\documentclass[11pt]{article}
\usepackage[T1]{fontenc}
\usepackage{textcomp}
\usepackage[margin=0.75in]{geometry}
\usepackage{amsmath,amssymb,amsthm}
\usepackage{array,booktabs}
\usepackage{hyperref}
\setlength{\parindent}{0pt}
\newtheorem{theorem}{Theorem}
\theoremstyle{definition}
\newtheorem{definition}{Definition}
\title{Flow Proof Artifact: prop-12-perpendicular-from-point}
\author{Generated by \texttt{flow doc proof}}
\date{Flow Proof Book}
\begin{document}
\maketitle
To draw a perpendicular to a given infinite straight line from a given point not on it.\\[0.5em]
\textbf{Source.} euclid — Elements, Book I, Proposition 12\\[0.5em]
\bigskip
\noindent{\large \textbf{Derived fact 1.} \textit{To draw a perpendicular to a given infinite straight line from a given point not on it} \hfill the Euclidean plane \cdot Euclid Book I \cdot Proposition 12: drop a perpendicular from a point to a line}\par
\smallskip\noindent\textit{Coordinate: the Euclidean plane \cdot Euclid Book I \cdot Proposition 12: drop a perpendicular from a point to a line}\par
\smallskip\noindent\textit{Source: euclid — Elements, Book I, Proposition 12}\par
\smallskip\noindent\textit{Built on: proposition 2: a segment equal to a given segment, for Book I of the Elements in on the Euclidean plane, proposition 10: bisect a given segment, for Book I of the Elements in on the Euclidean plane, proposition 11: erect a perpendicular at a point on a line, for Book I of the Elements in on the Euclidean plane}\par
\medskip
\noindent\textbf{Goal.} To draw a perpendicular to a given infinite straight line from a given point not on it\par
\noindent$\displaystyle \angle CEB = 90^\circ$\par
\medskip
\noindent
\begin{tabular}{@{} >{\raggedright\arraybackslash}p{0.06\textwidth} >{\raggedright\arraybackslash}p{0.47\textwidth} >{\raggedleft\arraybackslash}p{0.41\textwidth} @{}}
\textbf{\#} & \textbf{Proof} & \textbf{Mathematics} \\
\midrule
\textbf{1.} & We prove that proposition 12: drop a perpendicular from a point to a line for Book I of the Elements in the Euclidean plane. & \\[0.45em]
\textbf{2.} & Let circle\_center\_C\_radius\_CD = circle\_at\_C\_with\_radius\_CD. & \\[0.45em]
\textbf{3.} & Let E = second\_intersection\_of\_circle\_with\_line\_AB. & \\[0.45em]
\textbf{4.} & Let F = midpoint\_of\_EF\_by\_Proposition\_10. & \\[0.45em]
\textbf{5.} & Let perpendicular\_at\_F = line\_through\_F\_perpendicular\_to\_AB\_by\_Proposition\_11. & \\[0.45em]
\textbf{6.} & We invoke the axiom governing Book I of the Elements in the Euclidean plane: proposition 2: a segment equal to a given segment, for Book I of the Elements in on the Euclidean plane. & \\[0.45em]
\textbf{7.} & We invoke the derived fact governing Book I of the Elements in the Euclidean plane: proposition 10: bisect a given segment, for Book I of the Elements in on the Euclidean plane. & \\[0.45em]
\textbf{8.} & We invoke the derived fact governing Book I of the Elements in the Euclidean plane: proposition 11: erect a perpendicular at a point on a line, for Book I of the Elements in on the Euclidean plane. & \\[0.45em]
\textbf{9.} & From step 2, step 3, step 4, step 5, step 6, step 7, and step 8, by the chain of results established in Book I (step 2, step 3, step 4, step 5, step 6, step 7, and step 8), we obtain angle CEB equals one right angle. Hence proven. & $\displaystyle \angle CEB = 90^\circ$ \\[0.45em]
\bottomrule
\end{tabular}
\label{thm:Geometry-Euclid Book I-proposition_12_drop_a_perpendicular_from_a_point_to_a_line}
\par\medskip\hrule
\end{document}
